\chapter{pysisyphus run report (menu 33)}
\label{ch:menu33}
\menulabel{33}

\section{Purpose}

pysisyphus writes no single result file: the record of a run is its console
(the cycle table, the outcome marker, the closing summary) plus a small set of
artifacts next to it -- the trajectory, the closing geometry, the normalised
run record \file{RUN.yaml} and the per-call engine outputs under
\file{qm\_calcs/}.  This menu reads those back and cross-checks them against
each other, so that a run can be brought into the toolkit's chains (the
closing geometry goes on to any structure-reading menu; the per-call ORCA
outputs are ordinary ORCA outputs for the menu-1 check-up).

\section{What you need}

The run directory, or the console capture alone (e.g.\ \code{pysis x.yaml >
run.out}).  Given a directory, the menu finds the capture itself (the file
carrying the \code{\# RUNNING ... \#} banner), then picks up the trajectory,
\file{RUN.yaml}, the closing geometry and the calculator outputs when they are
present.

\section{Walkthrough}

\lstinputlisting[style=fbkcode, firstline=39, lastline=69,
  title={The menu-33 example (the H$_2$O optimisation: the run's own facts,
  the cycle table, the outcome, the closing state and the five cross-checks)}]
  {examples/expected/m33_pysisyphus.txt}

\section{How to read it}

\begin{itemize}
  \item the outcome is the program's own marker -- \code{Converged!} or
    \code{Number of cycles exceeded!}.  It is deliberately not derived from
    the force values: a run stopped by its cycle limit can close with forces
    already below the thresholds and is still not converged (measured on the
    stopping fixture);
  \item a stopped run evaluates \emph{one extra geometry} (the closing state)
    after its last cycle: the console has one cycle fewer than there are
    calculator calls, and the closing geometry file belongs to that extra
    evaluation, not to any trajectory frame (the menu says so instead of
    pretending a match);
  \item the cross-checks tie the artifacts together: the trajectory has one
    frame per cycle row; the closing frame's energy equals the calculator call
    that produced it; the \code{Final summary} equals the run's last call; and
    the closing geometry file reproduces the closing state.  The calculator
    calls are the ORCA outputs under \file{qm\_calcs/}, read by this toolkit's
    ORCA parser -- the run's energies are re-derived from ORCA's files rather
    than trusted to the console text.
\end{itemize}

\section{Crashed runs and the boundary survey (Wave 4.6)}

\begin{readbox}
\begin{itemize}
  \item a crashed run is classified: the measured ORCA-6 Hessian parse stop is
    recognised by its \code{\$multiplicity} signature (the \file{crashed\_*}
    backup directory is searched for the offending \file{.hess}), reported
    together with the way forward (a model Hessian; frequency checks outside
    pysisyphus); other crashes point at the backup directory copied next to
    the capture;
  \item when the crash followed a converged optimisation, the report says so:
    the geometry part finished, the post-processing did not;
  \item \file{optimization.h5} carries the full per-cycle history but is
    HDF5, which is not a dependency of this toolkit; the reader takes the same
    numbers from the text artifacts and names the file as unread;
  \item \file{qm\_calcs/cur\_out} is a symlink into the scratch directory,
    which pysisyphus cleans after every call -- it dangles after the run; the
    per-call copies (\file{calculator\_*.orca.*}) are the durable ones.
\end{itemize}
\end{readbox}
